\documentclass[letterpaper,11pt,openany,oneside]{report}	% informe estandar
%\documentclass[letterpaper,12pt,openany,oneside,notitlepage]{report}	% sin portada

\usepackage{dsfont}
\usepackage[utf8x]{inputenc}							% codificacion español
\usepackage[activeacute,spanish]{babel}								% lenguage español
\usepackage{latexsym}											% simbolos
\usepackage{amssymb}											% simbolos AMS
\usepackage{float}												% objetos flotantes
\usepackage[center,footnotesize]{caption}	% titulos de figuras y tablas[posicion,tamaño]
\usepackage{graphicx}							% soporte figuras
\usepackage{epstopdf}
%\usepackage[pdftex]{graphicx}							% soporte figuras
\usepackage[nottoc,numbib]{tocbibind}			% bibliografia en indice
\usepackage{sectsty}											% soporte estilizar titulos
\usepackage[small,compact]{titlesec}			% soporte cambiar tamaño titulos [tamaño letra, espaciamiento]
%\usepackage{pdfpages}						% soporte anexar pdf externo con \includepdf{archivo.jpg}
%\usepackage{savetrees}										% menos papel
\usepackage{lscape}
\usepackage{subfig}
\usepackage[pdftex,colorlinks,bookmarksnumbered,bookmarksopen=true,plainpages=false,linkcolor=black,citecolor=red,urlcolor=blue]{hyperref}
\usepackage{amsmath}
\usepackage{amstext}
\usepackage{textcomp}
\usepackage{mathtools}
\usepackage{psfrag}
\usepackage{amsfonts}
\usepackage{bbm}

% soporte enlaces, colores default: linkcolor=red, citecolor=green, urlcolor=magenta

\newcommand{\be}{\begin{equation}}	% comenzar ecuacion con '\be'
\newcommand{\ee}{\end{equation}}		% finalizar ecuacion con '\ee'
\newcommand{\itemcheck}{\item[\checkmark]} 	%crear un tick para los items
\newcommand{\graficocaption}[4]{
\begin{figure}[htbp]
\centering
\includegraphics[height=#1]{#2}
\caption{#3}
\label{#4}
\end{figure}
}
\newcommand{\graficocaptiont}[4]{
\begin{figure}[htbp]
\centering
\includegraphics[height=#1, trim=#4, clip]{#2}
\caption{#3}
\end{figure}
}

\voffset -1.2in
\oddsidemargin 0in
\textheight 8.5in
\textwidth 6.5in
\linespread {1.3}

\titleformat{\chapter}[display] % cambiamos el formato de los capitulos
{\bfseries\Large} % por defecto se usaron caracteres de tamaño \Large en negrita
{% contenido de la etiqueta
 \titlerule % linea horizontal
 %\filleft % texto alineado a la derecha
 \filright % texto alineado a la izquierda
 \large\chaptertitlename\ % "Capitulo" o "Apendice" en tamaño \normalsize en lugar de \Large
 \large\thechapter} % numero de caputulo en tamaño \normalsize
{0mm} % espacio minimo entre etiqueta y cuerpo
%{\filleft} % texto del cuerpo alineado a la derecha
{\filright\Large} % texto del cuerpo alineado a la izquierda
[\vspace{0.5mm} \titlerule] % despues del cuerpo, dejar espacio vertical y trazar linea horizontal gruesa

\captionsetup{labelfont=sf,textfont=sf,justification=justified,format=hang}



%simbolos %%Matias Courdurier%%

\newcommand{\R}{\mathbb{R}}		% Reales
\newcommand{\N}{\mathbb{N}}		% Naturales
\newcommand{\one}{\mathbbm{1}}		% Naturales


\begin{document}

\renewcommand{\contentsname}{\'Indice}	% cambia nombre: 'Indice general' -> 'Indice'
\renewcommand{\tablename}{Tabla}			% cambia nombre: 'Cuadro' -> 'Tabla'
\allsectionsfont{\sffamily}	% titulos en sans-serif, necesita paquete sectsty
\sffamily										% texto en sans-serif


%%%%%%%%%%%%%%%%%%%%%%%PORTADA%%%%%%%%%%%%%%%%%%%%%%%
\title{
\vskip -2cm
\includegraphics[height=1.5cm]{portada/figuras/logo_mmgeo}\hspace*{1.2cm}
\includegraphics[width=5cm]{./portada/figuras/logo_cmm} \hfill
\includegraphics[height=2cm]{./portada//figuras/logo_dim}\\
{\small Av. Blanco Encalada 2120 \hfill  \, \\
Santiago, Chile \hfill \, \\
Fono: 978 4525  \hfill \, \\
Fax: 688 9705  \hfill \, \\ }	
	\vspace{0.5in}
	\vspace{0.2in}
	\hrule
	\vspace{0.2in}
	\Huge Un estudio para la determinaci\'on de la superficie de quiebre en una mina subterr\'anea en El Teniente
	\vspace{0.2in}
	\hrule
	\vspace{0.2in}
	\LARGE Informe Final
	\vspace{\stretch{1}}
}

\author{
	\large Ra\'ul Man\'asevich  (director)\\	
    \large Roger Bustamante\\
    \large Javier Ruiz\\
    \large Duvan Henao\\    
    \large Matias Courdurier\\
    \large Claudia Vallejos \\
    \large Matias Valdenegro \\
    \large Lenardo (Apoyo a Duvan ....) \\
    \large Lenardo (Apoyo a Matias Courdurier ......) \\
    \large Mael (apoyo de Duvan.......) \\    
    \large Cristian Jara\\
	\large Mario Fern\'andez\\
	%\vspace{\stretch{1}}
}

\date{\large\textbf{\today}}
%%%%%%%%%%%%%%%%%%%%%%%%%%%%%%%%%%%%%%%%%%%%%
%%%%%%%%%%%%%%%%%%%%%%%%%%%%%%%%%%%%%%%%%%%%%

\maketitle
\tableofcontents
\thispagestyle{empty}	% correr primera pagina despues del indice, funciona solo si el indice tiene una pagina
\setcounter{page}{1}

%%%%%%%%%%%%%%%%Documento%%%%%%%%%%%%%%%%%%%%%%%%%%%%%%%%%%%%%%%%%%%%%%%%%%%%%%%%%%%%%%%%%%%%%%%


%% Raul Manasevich
\chapter{Resumen Ejecutivo}
El presente informe da cuenta de los resultados del proyecto \emph{Un estudio para la determinaci\'on de la superficie de quiebre en una mina subterr\'anea en El Teniente}, realizado por acad\'emicos e ingenieros de la Facultad de Ciencias F\'isicas y Matem\'aticas de la Universidad de Chile, agrupados en torno al Laboratorio de Modelamiento Matem\'atico para Geomec\'anica (MMGeo) perteneciente al Centro de Modelamiento Matem\'atico (CMM), y con el apoyo en recursos humanos y financieros de la Superintendencia de Geomec\'anica de la Divisi\'on El Teniente, Codelco. \\
El proyecto tiene como objetivo la evaluaci\'on de factibilidad de distintos m\'etodos de estimaci\'on de la superficie de quiebre. Estos m\'etodos son:






%% Duvan Henao
%% Responsable : Matias Courderier 

\chapter{Linea Time Reverse}

\section{Resumen}

\section{El Modelo}

\section{Descripci\'on Software}


\section{Casos de estudio}

\subsection{Caso Pilar Norte}


El primer ejemplo que se presenta, consiste en ...


\subsubsection{Input}

\subsubsection{Output}

\subsubsection{Visualizaci\'on}

\subsubsection{An\'alisis de los resultados}

\subsection{Caso Esemeralda}


El primer ejemplo que se presenta, consiste en ...


\subsubsection{Input}

\subsubsection{Output}

\subsubsection{Visualizaci\'on}

\subsubsection{An\'alisis de los resultados}
% \input{cme3d_comp.tex} 


\section{Resultados Generales}

\section{Extensiones futura}

\section{Manual de Uso}

%% Roger Bustamante
%% Responsable : Matias Courderier 

\chapter{Linea Time Reverse}

\section{Resumen}

\section{El Modelo}

\section{Descripci\'on Software}


\section{Casos de estudio}

\subsection{Caso Pilar Norte}


El primer ejemplo que se presenta, consiste en ...


\subsubsection{Input}

\subsubsection{Output}

\subsubsection{Visualizaci\'on}

\subsubsection{An\'alisis de los resultados}

\subsection{Caso Esemeralda}


El primer ejemplo que se presenta, consiste en ...


\subsubsection{Input}

\subsubsection{Output}

\subsubsection{Visualizaci\'on}

\subsubsection{An\'alisis de los resultados}
% \input{cme3d_comp.tex} 


\section{Resultados Generales}

\section{Extensiones futura}

\section{Manual de Uso}

%% Matias Courdurier
%% Responsable : Matias Courderier 

\chapter{Linea Time Reverse}

\section{Resumen}

\section{El Modelo}

\section{Descripci\'on Software}


\section{Casos de estudio}

\subsection{Caso Pilar Norte}


El primer ejemplo que se presenta, consiste en ...


\subsubsection{Input}

\subsubsection{Output}

\subsubsection{Visualizaci\'on}

\subsubsection{An\'alisis de los resultados}

\subsection{Caso Esemeralda}


El primer ejemplo que se presenta, consiste en ...


\subsubsection{Input}

\subsubsection{Output}

\subsubsection{Visualizaci\'on}

\subsubsection{An\'alisis de los resultados}
% \input{cme3d_comp.tex} 


\section{Resultados Generales}

\section{Extensiones futura}

\section{Manual de Uso}

%% Javier Ruiz
%% Responsable : Matias Courderier 

\chapter{Linea Time Reverse}

\section{Resumen}

\section{El Modelo}

\section{Descripci\'on Software}


\section{Casos de estudio}

\subsection{Caso Pilar Norte}


El primer ejemplo que se presenta, consiste en ...


\subsubsection{Input}

\subsubsection{Output}

\subsubsection{Visualizaci\'on}

\subsubsection{An\'alisis de los resultados}

\subsection{Caso Esemeralda}


El primer ejemplo que se presenta, consiste en ...


\subsubsection{Input}

\subsubsection{Output}

\subsubsection{Visualizaci\'on}

\subsubsection{An\'alisis de los resultados}
% \input{cme3d_comp.tex} 


\section{Resultados Generales}

\section{Extensiones futura}

\section{Manual de Uso}

%% Cludia Vallejos
%% Responsable : Matias Courderier 

\chapter{Linea Time Reverse}

\section{Resumen}

\section{El Modelo}

\section{Descripci\'on Software}


\section{Casos de estudio}

\subsection{Caso Pilar Norte}


El primer ejemplo que se presenta, consiste en ...


\subsubsection{Input}

\subsubsection{Output}

\subsubsection{Visualizaci\'on}

\subsubsection{An\'alisis de los resultados}

\subsection{Caso Esemeralda}


El primer ejemplo que se presenta, consiste en ...


\subsubsection{Input}

\subsubsection{Output}

\subsubsection{Visualizaci\'on}

\subsubsection{An\'alisis de los resultados}
% \input{cme3d_comp.tex} 


\section{Resultados Generales}

\section{Extensiones futura}

\section{Manual de Uso}

%% Mario Fernandez
%% Responsable : Matias Courderier 

\chapter{Linea Time Reverse}

\section{Resumen}

\section{El Modelo}

\section{Descripci\'on Software}


\section{Casos de estudio}

\subsection{Caso Pilar Norte}


El primer ejemplo que se presenta, consiste en ...


\subsubsection{Input}

\subsubsection{Output}

\subsubsection{Visualizaci\'on}

\subsubsection{An\'alisis de los resultados}

\subsection{Caso Esemeralda}


El primer ejemplo que se presenta, consiste en ...


\subsubsection{Input}

\subsubsection{Output}

\subsubsection{Visualizaci\'on}

\subsubsection{An\'alisis de los resultados}
% \input{cme3d_comp.tex} 


\section{Resultados Generales}

\section{Extensiones futura}

\section{Manual de Uso}


%% Referecias

\begin{thebibliography}{99}

%% ejemplos...
	\bibitem{detcav} Superintendencia de Geomec\'anica TTE-GRMD. Modelo para monitoreo de Caving. Noviembre 2010. 
	
	\bibitem{ppa} MMGeo. Estudio 2D de Escenarios de Inter\'es en el Pre-acondicionamiento y Determinaci\'on de la Cavidad M\'axima Estable en 3D. Marzo, 2011.
	\bibitem{cundall} P.A. Cundall, O.D. Strack, A discrete numerical model for granular assemblies, Geotechnique 29 (1979) 47-65.

  \bibitem{PrOrVa} A.I. Prilepko, D.G. Orlovskii, I.A. Vasin, {\em Methods for Solving Inverse Problems in Mathematical
Physics}, Marcel Dekker, New York, 2000.

%%%%%  Fin Ejemplos %%%%%

%% Linea Duvan

%% Linea Roger

%% Linea Matias

%% linea Javier

%% linea Claudia

%% linea Mario


\end{thebibliography}

\end{document}
